% answers_short.tex —— 2004 年·普通高等学校招生考试·重庆卷·理科 · 简短答案速查
% 由 tools/make_short_answers.py 自 answers.tex 的 \gksolhead/\gkanswer 提取，
% 并与 2004/2004_数学真题_重庆理科_答案要点.md 逐题比对后生成，勿手改（改源 MD 后重跑）。

\begin{gknotice}
\textbf{使用说明}：本节只给最终答案便于快速核对，完整推导见解析版。答案为编写者独立推导，\textbf{非官方标准答案}。
\end{gknotice}

\gkpart{选择题}
\begin{center}\begin{tabular}{ccccccccccccc}
\toprule
题号 & 1 & 2 & 3 & 4 & 5 & 6 & 7 & 8 & 9 & 10 & 11 & 12 \\
\midrule
答案 & D & A & D & A & B & C & C & C & B & B & B & D \\
\bottomrule
\end{tabular}\end{center}

\gkpart{填空题}
\begin{itemize}[leftmargin=2.6em]
  \item[\textbf{13.}] $-2$
  \item[\textbf{14.}] $\dfrac{\pi}{4}$
  \item[\textbf{15.}] $\dfrac{\pi}{3}$
  \item[\textbf{16.}] $[-1, 3]$
\end{itemize}

\gkpart{解答题}
\begin{itemize}[leftmargin=2.6em]
  \item[\textbf{17.}] 最小正周期为 $\pi$，最小值为 $-2$；单调递增区间为 $\left[0, \dfrac{\pi}{3}\right]$ 和 $\left[\dfrac{5\pi}{6}, \pi\right]$。
  \item[\textbf{18.}] (1) 分布列见解析，期望 $E\xi = \dfrac{525}{256}$；(2) $\dfrac{175}{256}$
  \item[\textbf{19.}] (1) 见解析版；(2) $\dfrac{\sqrt{5}}{10}$
  \item[\textbf{20.}] (1) 见解析版；(2) $[2, +\infty)$
  \item[\textbf{21.}] 见解析证明；圆 $H$ 的面积最小时直线 $AB$ 的方程为 $x = 2p$。
  \item[\textbf{22.}] (1) 见解析版；(2) $b_n > b_{n+1}$（数列 $\{b_n\}$ 单调递减）
\end{itemize}
